\dish{ラーメン}
\altdish{拉麺}
\altdish{Ramen, noodles (homemade)}
\pour{4}
\prep{2~h + repos (1--2~jours) + cuisson (90 secondes)}
\source{cjm}
\label{ラーメン}

\begin{ingredients}
  \ingr{4}{g}{sel}
  \ingr{175}{ml}{eau}
  \ingr{350}{g}{farine T55}
  \ingr{50}{g}{fécule de maïs (ou arrow-root ou tapioca)}
  \ingrS{4}{g}{bicarbonate de soude (à défaut de kansui, voir ci-dessous)}
  \ingr{}{}{fécule de pomme de terre ou de maïs pour le façonnage}
\end{ingredients}


\begin{recipe}
  \begin{enumerate}

  \item Dissoudre le sel dans l'eau.  Ajouter la farine, la fécule et
    le bicarbonate.  Former une boule et laisser reposer dans
    un torchon humide.

  \item Mettre dans un sac en plastique et pétrir avec les pieds pour
    obtenir une pâte lisse.  Laisser reposer 10~minutes, puis répéter
    cette opération 4~fois.

  \item Laisser reposer 1~heure en été, 3~heures en hiver.

  \item Rouler à 1~cm d'épaisseur.  Fariner de fécule.  Rouler comme
    des pâtes italiennes.  Former des spaghettis plutôt épais (3--4
    sur ma machine à pâte).

  \item Laisser reposer 1--2~jours dans un sac en plastique.

  \item Faire cuire 90~secondes dans de l'eau bouillante salée.
    Égoutter et servir.

  \end{enumerate}
\end{recipe}

\stage{Kansui (かん水 / 鹹水)}\nopagebreak

Solution alcaline de carbonates de sodium et de potassium.  Elle
porte le pH de la pâte à 9--11 : le gluten se raffermit (pâtes
élastiques qui tiennent à la cuisson), la farine jaunit sans œufs, et
le goût devient celui des ramen.  Le bicarbonate de soude est moins
alcalin, mais cuit il devient du carbonate de sodium :
\[
  2\,\mathrm{NaHCO_3} \longrightarrow \mathrm{Na_2CO_3} + \mathrm{CO_2} + \mathrm{H_2O}
\]

\begin{ingredients}
  \ingr{100}{g}{bicarbonate de soude}
\end{ingredients}

\begin{recipe}
  \begin{enumerate}

  \item Étaler sur une plaque.  Cuire 1~heure à 120--150\degreeC.  Le
    poids doit tomber à environ 63~g.

  \item Conserver dans un bocal hermétique.

  \end{enumerate}
\end{recipe}

Utiliser 1~\% du poids de farine (4~g ici).  Plus caustique que le
bicarbonate : éviter le contact avec la peau et les yeux.

%\accord{}
